% ==============================================================================
% Fichier     : Contenu_PIU.tex

% ==============================================================================

\section{Introduction}

Le \textbf{Processus Investigatif Universel (PIU)} est une architecture méthodologique modulaire garantissant la rigueur de toute enquête (organisationnelle, numérique ou judiciaire). Il structure l'investigation comme un processus scientifique reproductible reposant sur quatre piliers : \textbf{objectivité factuelle}, \textbf{intégrité probatoire}, \textbf{traçabilité complète} et \textbf{défendabilité}. Selon les enjeux du cas, l'enquêteur choisit un niveau de profondeur : \textit{Minimal} (incidents internes mineurs), \textit{Standard} (fraudes ou incidents de sécurité) ou \textit{Avancé} (affaires criminelles à fort impact).
\section{Les Quality Gates (QG)}

Les \textbf{Quality Gates} sont des points de contrôle validant la transition entre chaque phase. Leur évaluation est tri-dimensionnelle : validité pré-opérationnelle, l'intégrité probatoire et la robustesse analytique. Un système de feux guide la décision — \textbf{Vert} : poursuite, \textbf{Orange} : corrections requises, \textbf{Rouge} : arrêt ou retour en arrière.

\begin{qgbox}{Mécanismes de Contrôle}
    \begin{itemize}
        \item \textbf{Validité pré-opérationnelle:} Les protocoles de collecte ont-ils été respectés ?
        \item \textbf{Intégrité probatoire:} Les preuves supportent-elles réellement les hypothèses ?
        \item \textbf{Robustesse analytique :} A-t-on cherché activement à infirmer nos théories ?
    \end{itemize}
\end{qgbox}

\section{Phases Opérationnelles}
Avant tout, il faut préciser que tout processus investigatif commence avec un fait. 
\subsection{Phase I — Initialisation et Cadrage}
Évaluation de la recevabilité, signature d'un mandat définissant le périmètre (SLA) et planification des ressources. Le \textbf{QG1} valide le droit d'agir et les moyens disponibles.

\subsection{Phase II — Acquisition et Préservation}
Application de la règle \textbf{3-2-1} (3 copies, 2 supports, 1 hors site) et hachage cryptographique de chaque preuve. Le \textbf{QG2} verrouille la chaîne de conservation (\textit{Chain of Custody}).

\subsection{Phase III — Investigation Analytique}
Cœur intellectuel du PIU : entretiens structurés (\textbf{PEACE}), Analyse des Hypothèses Concurrentes (\textbf{ACH}) et \textbf{falsification active} (principe de Popper — réfuter l'hypothèse pour la valider). Le \textbf{QG3} valide la robustesse par \textit{Red Teaming}.

\subsection{Phase IV — Formalisation et Transmission}
Rapport structuré en trois volumes : synthèse décisionnelle, analyse détaillée et inventaire des preuves. Transmission obligatoire via canaux chiffrés.

\subsection{Phase V — Suivi et Amélioration Continue}
Maintien des preuves pour la phase contradictoire et \textbf{REX} (Retour d'Expérience) pour corriger les défaillances du processus.

\section{Conclusion}

Le PIU transfigure totalement le processus investigatif et permet de produire des résultats falsifiables et recevables pour la justice. Il est parsemé de verrous de contrÔle assurant la bonne conduite de l'investigation et en cas d'erreur permet de se retrouver sur ses pas. 